\documentclass[aspectratio=169, 10pt]{beamer}
\usetheme{metropolis}

% --- Edinburgh colours ----------------------------------------
\definecolor{UoERed}{HTML}{7A2318}
\definecolor{UoEGold}{HTML}{B8860B}
\definecolor{CodeBG}{HTML}{F7F3ED}
\definecolor{CodeFrame}{HTML}{D5C9B1}

\setbeamercolor{palette primary}{bg=UoERed, fg=white}
\setbeamercolor{frametitle}{bg=UoERed, fg=white}
\setbeamercolor{title separator}{fg=UoEGold}
\setbeamercolor{progress bar}{fg=UoEGold, bg=UoEGold!20}
\setbeamercolor{alerted text}{fg=UoERed}
\setbeamercolor{block title}{bg=UoERed!15, fg=UoERed}
\setbeamercolor{block body}{bg=UoERed!5}

% --- Packages -------------------------------------------------
\usepackage[T1]{fontenc}
\usepackage{lmodern}
\usepackage{amsmath, amssymb, mathtools}
\usepackage{listings}
\usepackage{graphicx, booktabs}
\usepackage{tikz}
\usetikzlibrary{arrows.meta, positioning, calc, decorations.pathreplacing}
\usepackage{hyperref}
\hypersetup{colorlinks=true, linkcolor=UoERed, urlcolor=UoERed!70!black}
\setbeamertemplate{navigation symbols}{}

\lstset{
  language=Python,
  basicstyle=\ttfamily\scriptsize,
  keywordstyle=\color{UoERed}\bfseries,
  commentstyle=\color{gray}\itshape,
  stringstyle=\color{UoEGold},
  backgroundcolor=\color{CodeBG},
  frame=single,
  rulecolor=\color{CodeFrame},
  numbers=none, breaklines=true, showstringspaces=false, tabsize=4,
  xleftmargin=4pt, xrightmargin=4pt, aboveskip=6pt, belowskip=4pt,
}

\AtBeginSection[]{
  \begin{frame}
    \vfill\centering
    \usebeamerfont{title}\insertsectionhead\par
    \vfill
  \end{frame}
}

\title{Week 3 --- Unconstrained Optimisation II}
\subtitle{Gradient Descent \& Newton's Method}
\author{Dr Juan Zurita}
\institute{Optimisation \& Mathematical Methods for Economics\\
The University of Edinburgh~$\cdot$~School of Economics}
\date{}

\begin{document}
\maketitle

\begin{frame}{Gradient Descent}
\textbf{Algorithm:} $\mathbf{x}_{k+1} = \mathbf{x}_k - \alpha\,\nabla f(\mathbf{x}_k)$\\\bigskip\begin{itemize}\item $\alpha > 0$: step size (learning rate)\item Follows the steepest descent direction\item \textbf{Linear} convergence rate\item Simple but can be slow on ill-conditioned problems\end{itemize}
\end{frame}

\begin{frame}{Newton's Method}
Uses second-order information:\\ $\mathbf{x}_{k+1} = \mathbf{x}_k - [H_f(\mathbf{x}_k)]^{-1}\nabla f(\mathbf{x}_k)$\\\bigskip\begin{itemize}\item \textbf{Quadratic} convergence near optimum\item Each step is more expensive (Hessian computation)\item May diverge if started far from optimum\end{itemize}
\end{frame}

\begin{frame}{Quasi-Newton: BFGS}
\textbf{Idea:} approximate the inverse Hessian from gradient differences.\\\bigskip\begin{itemize}\item No need to compute or invert the Hessian\item \textbf{Superlinear} convergence\item L-BFGS: limited-memory variant for large problems\item The workhorse of smooth optimisation in practice\end{itemize}\par\bigskip In Python: \texttt{scipy.optimize.minimize(method='BFGS')}
\end{frame}

\begin{frame}{Convergence Comparison}
\begin{center}\begin{tabular}{lcc}\toprule Method & Convergence & Cost per step \\ \midrule Gradient descent & Linear & $O(n)$ \\ Newton & Quadratic & $O(n^3)$ \\ BFGS & Superlinear & $O(n^2)$ \\ \bottomrule\end{tabular}\end{center}\par\bigskip Choose based on problem size and smoothness.
\end{frame}

\begin{frame}{Summary}
\begin{itemize}\item GD: simple, slow; Newton: fast, expensive; BFGS: best of both\item Step size matters --- too large diverges, too small crawls\item \texttt{scipy.optimize} provides production-quality implementations\item \textbf{Next week:} constrained optimisation with Lagrange multipliers\end{itemize}
\end{frame}
\end{document}
